% 第15章第5节：正交约束下的矩阵逼近（正交逼近的最优解、旋转约束下的逼近）。

\section{正交约束下的矩阵逼近}
\label{sec:matrix-approximation-orthogonal}

前面的约束是秩预算，本节把约束换成正交性。给定方阵$\bm A\in\mathbb R^{d\times d}$，在全部正交矩阵$\bm R$（$\bm R^\top\bm R=\bm I$）中寻找离$\bm A$最近的一个，误差仍用Frobenius范数。正交矩阵保持长度，所以问题是：哪一个保长作用最接近给定的线性作用？例如测得的旋转因噪声不再正交时，就要把它修正为最近的正交矩阵。

正交矩阵组成的集合不是线性子空间，也不由秩描述，前两节的投影化归不能沿用。第~\ref{chap:polar-decomposition}章已经构造了极分解$\bm A=\bm W\bm H$的正交因子$\bm W$；本节证明它正是所求的最近正交矩阵，再讨论要求$\det\bm R=1$时答案怎样改变。

\subsection{正交逼近的最优解}
\label{sec:matrix-approximation-orthogonal-optimum}

\paragraph{主要思想}

先把距离化为一个更简单的目标。对任意正交矩阵$\bm R$，由$\lVert\bm R\rVert_{\mathrm F}^2=\operatorname{tr}(\bm R^\top\bm R)=\operatorname{tr}\bm I=d$，
\[
 \lVert\bm A-\bm R\rVert_{\mathrm F}^2
 =\lVert\bm A\rVert_{\mathrm F}^2+d-2\operatorname{tr}(\bm R^\top\bm A),
\]
前两项与$\bm R$无关，所以距离最小等价于在正交矩阵中最大化$\operatorname{tr}(\bm R^\top\bm A)$。这个目标关于$\bm R$是线性的；第~\ref{sec:ls-rank-k-trace}节选择子空间时的$\operatorname{tr}(\bm Y^\top\bm G\bm Y)$关于$\bm Y$是二次的，那里的迹极值界不能直接搬用，需要重新给出上界。

上界来自SVD。代入$\bm A=\bm U\bm\Sigma\bm V^\top$，令$\bm Z=\bm U^\top\bm R\bm V$。它是正交矩阵的乘积，仍然正交，且$\bm R=\bm U\bm Z\bm V^\top$；$\bm R$取遍全部正交矩阵时，$\bm Z$也取遍全部正交矩阵。于是问题变成在正交矩阵$\bm Z$中最大化$\operatorname{tr}(\bm Z^\top\bm\Sigma)=\sum_i\sigma_iz_{i,i}$。

\begin{theorem}[最佳正交逼近]
\label{thm:ls-best-orthogonal-approximation}
设$\bm A\in\mathbb R^{d\times d}$有SVD$\bm A=\bm U\bm\Sigma\bm V^\top$。对任意正交矩阵$\bm R$，
\[
 \operatorname{tr}(\bm R^\top\bm A)\leq\sum_{i=1}^d\sigma_i ,
\]
取$\bm R=\bm U\bm V^\top$时等号成立。因此
\[
 \min_{\bm R^\top\bm R=\bm I}\lVert\bm A-\bm R\rVert_{\mathrm F}^2
 =\sum_{i=1}^d(\sigma_i-1)^2 ,
\]
最小值由极分解的正交因子达到。$\bm A$可逆时最优解唯一，就是$\bm W=\bm A(\bm A^\top\bm A)^{-1/2}$。
\end{theorem}
\begin{proof}
令$\bm Z=\bm U^\top\bm R\bm V$。由迹的循环性，
\[
 \operatorname{tr}(\bm R^\top\bm A)
 =\operatorname{tr}(\bm V\bm Z^\top\bm U^\top\bm U\bm\Sigma\bm V^\top)
 =\operatorname{tr}(\bm Z^\top\bm\Sigma)
 =\sum_{i=1}^d\sigma_iz_{i,i}.
\]
$\bm Z$的每一列都是单位向量，所以$|z_{i,i}|\leq1$；又$\sigma_i\geq0$，故$\sigma_iz_{i,i}\leq\sigma_i$，求和即得上界。$\bm R=\bm U\bm V^\top$时$\bm Z=\bm I$，等号成立。

由式~\eqref{eq:matrix-analysis-frobenius-norm}，$\lVert\bm A\rVert_{\mathrm F}^2=\sum_i\sigma_i^2$，代入展开式得最小值
\[
 \sum_{i=1}^d\sigma_i^2+d-2\sum_{i=1}^d\sigma_i
 =\sum_{i=1}^d(\sigma_i-1)^2 .
\]
由定理~\ref{thm:ls-polar-general}，$\bm U\bm V^\top$是极分解的正交因子。

设$\bm A$可逆，则全部$\sigma_i>0$。等号要求每一项$\sigma_iz_{i,i}=\sigma_i$，即$z_{i,i}=1$；单位列向量$\bm Z_{:,i}$的第$i$个分量为$1$，其余分量只能为零，于是$\bm Z_{:,i}=\bm e_i$，$\bm Z=\bm I$，$\bm R=\bm U\bm V^\top$。由定理~\ref{thm:ls-polar-invertible}，它就是$\bm W$。
\end{proof}

去掉伸缩得到的正交因子，同时解决了“离$\bm A$最近的正交矩阵”这个问题；这一结论称为\term{最佳正交逼近}（best orthogonal approximation）。误差$\sum_i(\sigma_i-1)^2$逐个比较奇异值与$1$：正交矩阵的奇异值全为$1$，最佳正交逼近只是把$\bm\Sigma$换成$\bm I$。

$\bm A$奇异时，最优解一般不唯一，而且最优解恰好是全部极分解的正交因子。一方面，若$\bm A=\bm W'\bm H$是任一极分解，则$\operatorname{tr}(\bm W'^\top\bm A)=\operatorname{tr}(\bm W'^\top\bm W'\bm H)=\operatorname{tr}\bm H=\sum_i\sigma_i$，达到上界。

另一方面，设$r=\operatorname{rank}\bm A$，$\sigma_i=0$（$i>r$）的项对迹没有贡献，等号只要求$z_{i,i}=1$（$i\leq r$）。按上面的论证，$\bm Z$的前$r$列是$\bm e_1,\ldots,\bm e_r$；由正交性，其余各列与它们正交，于是$\bm Z=\operatorname{diag}(\bm I_r,\bm Z')$，其中$\bm Z'$是某个$d-r$阶正交矩阵。$\bm U\bm Z$仍正交，且只改变了$\bm U$的后$d-r$列，这些列在$\bm U\bm\Sigma$中都乘以零，所以$\bm A=(\bm U\bm Z)\bm\Sigma\bm V^\top$也是SVD，由定理~\ref{thm:ls-polar-general}，$\bm R=(\bm U\bm Z)\bm V^\top$是极分解的正交因子。不同的$\bm Z'$正对应第~\ref{chap:polar-decomposition}章所说零伸缩方向上不同的正交补全。

\paragraph{算法实现与算例}

计算只需三步：求$\bm A$的完整方阵SVD；直接相乘得到$\bm R_\star=\bm U\bm V^\top$；核验$\bm R_\star^\top\bm R_\star=\bm I$，并比较误差平方与$\sum_i(\sigma_i-1)^2$。这里没有限制$\det\bm R_\star$的符号，所得的可能是旋转，也可能包含反射。

\begin{example}[二阶正交逼近]
\label{ex:matrix-approximation-orthogonal-two-by-two}
取
\[
 \bm A=\begin{bmatrix}0&-2\\3&0\end{bmatrix}
 =\bm U\operatorname{diag}(3,2),\qquad
 \bm U=\begin{bmatrix}0&-1\\1&0\end{bmatrix},\quad\bm V=\bm I .
\]
逐项核对：$\bm U\operatorname{diag}(3,2)$的第一列是$3\bm U_{:,1}=(0,3)^\top$，第二列是$2\bm U_{:,2}=(-2,0)^\top$，正是$\bm A$；$\bm U$正交，奇异值$3\geq2$按降序排列。于是$\bm R_\star=\bm U\bm V^\top=\bm U$，残差
\[
 \bm A-\bm R_\star=\begin{bmatrix}0&-1\\2&0\end{bmatrix},
\]
平方误差为$1+4=5=(3-1)^2+(2-1)^2$。$\bm R_\star$是逆时针旋转$90$度，与$\det\bm A=6>0$一致。
\end{example}

第~\ref{chap:polar-decomposition}章贯穿使用的矩阵$\bm A=\left[\begin{smallmatrix}0&-1\\3&0\end{smallmatrix}\right]$也可以这样核对。那里$\bm S=\bm A^\top\bm A=\operatorname{diag}(9,1)$，奇异值为$3$与$1$，最小平方误差为$(3-1)^2+(1-1)^2=4$；直接计算也得$\bm A-\bm W=\left[\begin{smallmatrix}0&0\\2&0\end{smallmatrix}\right]$，其Frobenius范数的平方为$4$。含噪的旋转估计常用这一步，把一个近似正交的矩阵修正为正交矩阵。

\paragraph{算法分析}

除一次$d\times d$的SVD外，只需一次矩阵乘法，约$d^3$次乘法；SVD本身的代价取决于所用算法，本书不作分析。$\bm A$可逆时最优解唯一，与所取的SVD无关；$\bm A$奇异时，不同的SVD可能给出不同的最优正交矩阵，它们的差别只在零奇异值方向上，误差相同。

最优解是在全部正交矩阵中取得的。$\bm A$可逆时，唯一的最优解$\bm W$满足$\det\bm W$与$\det\bm A$同号（第~\ref{sec:ls-polar-orthogonal-factor}节），所以$\det\bm A<0$时它包含一次反射。若应用要求结果是旋转，就需要下一小节的修正。

\subsection{旋转约束下的逼近}
\label{sec:matrix-approximation-rotation}

\paragraph{主要思想}

刚体的姿态估计等问题要求结果是\emph{旋转}，即正交且$\det\bm R=1$；包含反射的正交矩阵不是允许的答案。于是在上一小节的迹问题中增加行列式约束：在$\bm R^\top\bm R=\bm I$、$\det\bm R=1$下最大化$\operatorname{tr}(\bm R^\top\bm A)$。

记$s=\det(\bm U\bm V^\top)=\det\bm U\det\bm V$。正交矩阵的行列式是$\pm1$，所以$s=\pm1$。若$s=1$，上一小节的$\bm U\bm V^\top$本身是旋转；它在更大的集合中最优，在旋转中当然也最优，行列式约束不起作用。若$s=-1$，$\bm U\bm V^\top$不可行，需要新的上界。这时对任一旋转$\bm R$，$\bm Z=\bm U^\top\bm R\bm V$正交，且
\[
 \det\bm Z=\det\bm U\cdot\det\bm R\cdot\det\bm V=s=-1 .
\]
上一小节的上界依赖$\bm Z=\bm I$可以取到，现在$\det\bm I=1$，这已不可能。下面的引理说明，行列式为$-1$的正交矩阵的迹至少比$d$小$2$。

\begin{lemma}[行列式为负的正交矩阵的迹]
\label{lem:matrix-approximation-reflection-trace}
设$\bm Z\in\mathbb R^{d\times d}$正交且$\det\bm Z=-1$，则存在单位向量$\bm x$使$\bm Z\bm x=-\bm x$，并且$\operatorname{tr}\bm Z\leq d-2$。
\end{lemma}
\begin{proof}
先证$-1$是特征值。由$\bm Z^\top\bm Z=\bm I$，$\bm Z+\bm I=\bm Z(\bm I+\bm Z^\top)=\bm Z(\bm I+\bm Z)^\top$，取行列式得
\[
 \det(\bm Z+\bm I)=\det\bm Z\cdot\det(\bm I+\bm Z)=-\det(\bm Z+\bm I),
\]
所以$\det(\bm Z+\bm I)=0$，存在单位向量$\bm x$使$\bm Z\bm x=-\bm x$。再由$\bm x=\bm Z^\top\bm Z\bm x=-\bm Z^\top\bm x$，得$\bm Z^\top\bm x=-\bm x$。

$d=1$时$\bm Z=(-1)$，$\operatorname{tr}\bm Z=-1=d-2$。设$d\geq2$，把$\bm x$扩充为$\mathbb R^d$的标准正交基，排成正交矩阵$\bm Q=[\bm x,\bm Q_2]$。由$\bm Z\bm x=-\bm x$与$\bm Z^\top\bm x=-\bm x$，$\bm Q_2^\top\bm Z\bm x=-\bm Q_2^\top\bm x=\bm0$，$\bm x^\top\bm Z\bm Q_2=(\bm Z^\top\bm x)^\top\bm Q_2=\bm0$，于是
\[
 \bm Q^\top\bm Z\bm Q=\begin{bmatrix}-1&\bm0\\\bm0&\bm Z_2\end{bmatrix},\qquad\bm Z_2=\bm Q_2^\top\bm Z\bm Q_2 .
\]
左端是正交矩阵的乘积，仍正交，所以$\bm Z_2$的各列是单位向量，对角元都不超过$1$，$\operatorname{tr}\bm Z_2\leq d-1$。迹在相似变换下不变，故$\operatorname{tr}\bm Z=-1+\operatorname{tr}\bm Z_2\leq d-2$。
\end{proof}

用特征值的语言看，同一结论是这样的：实正交矩阵的特征值模都为$1$，非实特征值成共轭对出现，每对的乘积为$1$、实部之和不超过$2$；实特征值只能是$\pm1$。行列式等于全部特征值之积，为$-1$时特征值$-1$必出现奇数次，至少一次；迹等于特征值之和，于是至多为$(d-1)\cdot1+(-1)=d-2$。上面的证明只用实向量完成了同一论证。

\begin{theorem}[最佳旋转逼近]
\label{thm:matrix-approximation-best-rotation}
设$\bm A\in\mathbb R^{d\times d}$有SVD$\bm A=\bm U\bm\Sigma\bm V^\top$，$\sigma_1\geq\cdots\geq\sigma_d\geq0$，$s=\det(\bm U\bm V^\top)$，$\bm D_s=\operatorname{diag}(1,\ldots,1,s)$。则对一切满足$\bm R^\top\bm R=\bm I$、$\det\bm R=1$的$\bm R$，
\[
 \operatorname{tr}(\bm R^\top\bm A)\leq\sum_{i=1}^{d-1}\sigma_i+s\,\sigma_d ,
\]
等号由$\bm R_\star=\bm U\bm D_s\bm V^\top$取到。因此
\[
 \min_{\bm R^\top\bm R=\bm I,\ \det\bm R=1}\lVert\bm A-\bm R\rVert_{\mathrm F}^2
 =\sum_{i=1}^{d-1}(\sigma_i-1)^2+(\sigma_d-s)^2 .
\]
\end{theorem}
\begin{proof}
$\bm R_\star$是正交矩阵的乘积，仍正交，且$\det\bm R_\star=\det\bm U\cdot s\cdot\det\bm V=s^2=1$，所以它是旋转。

$s=1$时$\bm D_s=\bm I$，上界就是定理~\ref{thm:ls-best-orthogonal-approximation}对全部正交矩阵给出的$\sum_i\sigma_i$，在旋转中自然成立，且由$\bm R_\star=\bm U\bm V^\top$取到。

设$s=-1$。对任一旋转$\bm R$，令$\bm Z=\bm U^\top\bm R\bm V$，则$\det\bm Z=-1$，由引理~\ref{lem:matrix-approximation-reflection-trace}，$\operatorname{tr}\bm Z\leq d-2$。与上一小节相同，$\operatorname{tr}(\bm R^\top\bm A)=\sum_i\sigma_iz_{i,i}$。由$|z_{i,i}|\leq1$得$1-z_{i,i}\geq0$，又$\sigma_i\geq\sigma_d$，所以
\[
 \sum_{i=1}^d\sigma_i-\operatorname{tr}(\bm R^\top\bm A)
 =\sum_{i=1}^d\sigma_i(1-z_{i,i})
 \geq\sigma_d\sum_{i=1}^d(1-z_{i,i})
 =\sigma_d(d-\operatorname{tr}\bm Z)\geq2\sigma_d ,
\]
即$\operatorname{tr}(\bm R^\top\bm A)\leq\sum_{i<d}\sigma_i-\sigma_d$。$\bm R=\bm R_\star$时$\bm Z=\bm D_s$，迹恰为$\sum_{i<d}\sigma_i-\sigma_d$，等号成立。

最后代入展开式。由$\lVert\bm A\rVert_{\mathrm F}^2=\sum_i\sigma_i^2$与$s^2=1$，
\[
 \sum_{i=1}^d\sigma_i^2+d-2\Bigl(\sum_{i=1}^{d-1}\sigma_i+s\sigma_d\Bigr)
 =\sum_{i=1}^{d-1}(\sigma_i-1)^2+(\sigma_d-s)^2 .
\]
\end{proof}

修正只发生在最小奇异方向上。$\bm R_\star=\bm U\bm D_s\bm V^\top$与无约束解$\bm U\bm V^\top$的差别，是把最后一对奇异方向的配对改为$\bm v_d\mapsto s\bm u_d$。若把符号改放在第$i$个方向上，同样得到旋转，但迹损失$2\sigma_i$；证明中的不等式说明，任何旋转的迹损失都至少是$2\sigma_d$，所以放在最小奇异方向上损失最少。

$s=-1$时，最小平方误差比无行列式约束时增加$(\sigma_d+1)^2-(\sigma_d-1)^2=4\sigma_d$。$\bm A$可逆时，$\det\bm A=\det\bm U\det\bm\Sigma\det\bm V$且$\det\bm\Sigma>0$，所以$s$就是$\det\bm A$的符号，与所取的SVD无关：$\det\bm A>0$时约束不起作用，$\det\bm A<0$时代价为$4\sigma_d$。

$\sigma_d=0$时，$s$依赖SVD的选取：把$\bm u_d$换成$-\bm u_d$不改变$\bm A$，却改变$\det\bm U$的符号。但这时$(\sigma_d-s)^2=1$与$s$无关，旋转约束的代价为零；而且两种选取给出同一个$\bm R_\star$，因为符号同时出现在$\bm u_d$与$s$中，乘积$s\bm u_d$不变。

唯一性分两种情形说明。$\bm A$可逆且$s=1$时，无约束问题的唯一解$\bm W$本身是旋转，所以它也是旋转约束下的唯一解。$\sigma_{d-1}>\sigma_d$时最优旋转也唯一：若$\sigma_d>0$且$s=-1$，上面不等式取等要求$(\sigma_i-\sigma_d)(1-z_{i,i})=0$，从而$i<d$时$z_{i,i}=1$，$\bm Z$的前$d-1$列是$\bm e_1,\ldots,\bm e_{d-1}$，最后一列只能是$\pm\bm e_d$，行列式为$-1$又迫使它为$-\bm e_d$，即$\bm Z=\bm D_s$；若$\sigma_d=0$，先取使$s=1$的SVD，取等要求$z_{i,i}=1$（$i<d$），同理$\bm Z=\operatorname{diag}(1,\ldots,1,\pm1)$，而$\det\bm Z=s=1$，故$\bm Z=\bm I$。

最小奇异值重复且$s=-1$，或$\operatorname{rank}\bm A\leq d-2$时，最优旋转可以不唯一，例子见下文。

\paragraph{算法实现与算例}

计算在上一小节的基础上增加两步。第一，求完整方阵SVD，奇异值按非增顺序排列，使最后一列$\bm u_d$、$\bm v_d$对应最小奇异值。第二，计算$s=\det\bm U\det\bm V$，构造$\bm D_s$，相乘得$\bm R_\star=\bm U\bm D_s\bm V^\top$。最后同时核验三项：$\bm R_\star^\top\bm R_\star=\bm I$，$\det\bm R_\star=1$，误差平方等于$\sum_{i<d}(\sigma_i-1)^2+(\sigma_d-s)^2$。

\begin{example}[对角矩阵的旋转逼近]
\label{ex:matrix-approximation-rotation-diagonal}
取$\bm A=\operatorname{diag}(3,-1)$，$\det\bm A=-3<0$。可选
\[
 \bm U=\operatorname{diag}(1,-1),\qquad\bm\Sigma=\operatorname{diag}(3,1),\qquad\bm V=\bm I ,
\]
核对$\bm U\bm\Sigma\bm V^\top=\operatorname{diag}(3,-1)=\bm A$。无行列式约束的最优正交矩阵是$\bm U\bm V^\top=\operatorname{diag}(1,-1)$，它是关于横轴的反射，误差平方为$\lVert\operatorname{diag}(2,0)\rVert_{\mathrm F}^2=4=(3-1)^2+(1-1)^2$。

$s=\det\bm U=-1$，取$\bm D_s=\operatorname{diag}(1,-1)$，得$\bm R_\star=\bm U\bm D_s\bm V^\top=\bm I$，误差平方为$\lVert\operatorname{diag}(2,-2)\rVert_{\mathrm F}^2=8$，比无约束时多$4=4\sigma_2$。直接比较全部二阶旋转：$\bm R=\left[\begin{smallmatrix}\cos\theta&-\sin\theta\\\sin\theta&\cos\theta\end{smallmatrix}\right]$给出$\operatorname{tr}(\bm R^\top\bm A)=3\cos\theta-\cos\theta=2\cos\theta$，只在$\theta=0$处取最大值$2=\sigma_1-\sigma_2$，所以$\bm R_\star=\bm I$是唯一的最优旋转，与$\sigma_1>\sigma_2$时的唯一性一致。
\end{example}

最小奇异值重复时，唯一性可能失败。取$\bm A=\operatorname{diag}(1,-1)$，奇异值为$1,1$，$s=-1$。对上面的任一旋转，$\operatorname{tr}(\bm R^\top\bm A)=\cos\theta-\cos\theta=0$，所有旋转的误差平方都是$\lVert\bm A\rVert_{\mathrm F}^2+2=4$，与公式$(1-1)^2+(1+1)^2=4$一致，但最优旋转有无穷多个。$\operatorname{rank}\bm A\leq d-2$时也有类似的自由度，例如$\bm A=\bm0$时每个旋转都最优。

\paragraph{算法分析}

与无约束情形相比，额外的工作只是求两个正交矩阵的行列式与一次对角缩放，主要代价仍是SVD。是否需要修正完全由$s$决定，不必比较多个候选。

最小奇异值接近零时，$s$对$\bm A$的微小变化很敏感：$\det\bm A$接近零，它的符号可能被扰动改变，最优旋转随之跳到另一个。但对同一个$\bm A$，两种符号给出的误差平方恰好相差$4\sigma_d$，这时也很小，最优误差本身随$\bm A$连续变化。所以这种情形下应报告最优误差，而不应把所得的旋转当作稳定的估计。
